% BLOCK SPECTRUM
\subsubsection{5.5 Complete equilibrium stability on the positive delay axis}
\label{complete-delay-spectrum}
The two fundamental crossings do not delimit the only possible stable
interval. Their recurrences can exchange stability again. Write
$P(\lambda)=(\lambda+a)(\lambda+d)(\lambda+c)$ and $k=C_ZB_E>0$,
where $a=-A_N$, $c=-C_E$ and $d=1/\tau_m$ are positive at the four
mobilising vectors. The characteristic equation is
$P(\lambda)-k\lambda e^{-\lambda\tau}=0$.

\textbf{Theorem 5.2 (Complete equilibrium root count).}
\emph{At each registered A/B, gated/ungated vector, let $\omega_-<\omega_+$
be the two positive frequencies and $\tau_-,\tau_+$ their fundamental
delays. Set $P_\pm=2\pi/\omega_\pm$, and let $n_\pm(\tau)$ count
$\tau_\pm+jP_\pm<\tau$, $j\ge0$. Away from these critical delays, the
number of characteristic roots in the open right half-plane, with
multiplicity, is exactly}
\[
 U(\tau)=2-2n_-(\tau)+2n_+(\tau).
\]
\emph{Every recurrence is simple; the lower-frequency family removes two
unstable roots and the higher-frequency family adds two. The complete
sets of exponentially stable positive delays are the following open
intervals (endpoints displayed approximately in years):}
\begin{center}\small
\begin{tabular}{lll}
\toprule Candidate & Gate & All stable intervals\\\midrule
A & gated & $(3.666149014,150.358477310)$; $(253.082596558,309.682275055)$\\
A & ungated & $(6.881410630,132.374903454)$; $(269.252385064,274.629634339)$\\
B & gated & $(5.512840731,80.424526714)$\\
B & ungated & $(6.213598734,76.290635688)$\\\bottomrule
\end{tabular}
\end{center}

\emph{Proof.} At zero delay the Routh array has two sign changes.
For $\operatorname{Re}\lambda\ge0$, $|\lambda|^3\le|P(\lambda)|
\le k|\lambda|$, so every such root lies in the fixed disk
$|\lambda|\le\sqrt{k}$; $\lambda=0$ is never a root. Thus roots cannot
enter from infinity as delay varies. The exhaustive modulus cubic
provides the only two imaginary frequencies. Implicit differentiation gives
\[
 \operatorname{Re}\bigl[(d\lambda/d\tau)^{-1}\bigr]
 =\frac{H'(\omega^2)}{|P(i\omega)|^2}.
\]
The interval signs are negative at $\omega_-$ and positive at
$\omega_+$, independently of recurrence. Counting crossings proves the
formula. Beyond
\[
 \tau_c=\frac{\tau_+/P_+-\tau_-/P_-}{1/P_+-1/P_-}
\]
the count is strictly positive. Interval ordering of the finitely many
crossings below $\tau_c$ proves the table; S17 gives the floor-count
bound and outward endpoint enclosures. \hfill$\square$

This closes the equilibrium spectral classification for all positive
delays, not merely a finite search range. In particular Candidate A has
a second, previously unexplored stability window. The five-regime
periodic-orbit description of Section 8 concerns the fundamental-delay
range only. The theorem does not count all Floquet multipliers of those
orbits or classify mobilising nonlinear basins.
% END SPECTRUM

% BLOCK BASIN
\subsubsection{7.5 A complete nonlinear protective basin}
\label{complete-protective-basin}
The stock-floor guarantee can be strengthened to an entire-basin
convergence theorem without assuming small perturbations. Write the
protective law as
\[
 \dot E=b(E)\{C(Z_{\rm meas})-E\},\qquad
 b(E)=\eta_p(1-E/M),\qquad
 C(Z)=\frac{E_0Z_{\rm ref}}{Z_{\rm ref}+Z},\quad M=E_{\max}.
\]
Here $Z_{\rm meas}=Z(t-\tau)$ for continuous deployment, or the last
reviewed delayed signal for native signal-hold. For frozen-effort review,
$E$ is held between reviews and the command is the exact nonlinear
flow of this scalar equation with $C$ frozen, not an Euler step.

\textbf{Theorem 7.5 (Complete protective equilibrium basin).}
\emph{Suppose the memory source $\Phi$ is nonnegative and globally
1-Lipschitz, with $\Phi(0)=\delta$, and put}
\[
 L_C=E_0/Z_{\rm ref},\qquad
 N_b=K(1-qE_0/r),\qquad
 \Gamma=qKL_C(1+K/N_b).
\]
\emph{If $0<E_0<M$, $qE_0<r$ and $\Gamma<1$, the unique positive equilibrium}
\[
 Z^*=\delta,\quad E^*=C(\delta),\quad N^*=K(1-qE^*/r)
\]
\emph{is globally asymptotically stable relative to exactly
$0<N(0)\le K$, $0\le E(0)<M$ and bounded nonnegative memory histories.
This holds for every positive adjustment rate and memory time, every
fixed finite deployment delay, and every positive fixed review cadence
for native signal-hold and exact nonlinear frozen-effort review. The
faces $N=0$ and $E=M$ are invariant and excluded from this basin.}

\emph{Proof outline.} The quota is bounded by $E_0$, so
$E\le\bar E=\max\{E(0),E_0\}<M$, $\limsup E\le E_0$ and
$\liminf N\ge N_b>0$. For $n=\log(N/N^*)$, $e=E-E^*$ and
$z=Z-\delta$, the stock equation is $\dot n=-a(t)n-qe$, with
$a(t)$ the logarithmic mean coefficient and asymptotic bounds
$rN_b/K\le a\le r$. If $L_x=\limsup|x(t)|$, the three input-output gains give
\[
 L_n\le\frac{qK}{rN_b}L_e,\qquad
 L_z\le K(rL_n+qL_e),\qquad L_e\le L_CL_z.
\]
Hence $L_e\le\Gamma L_e$, forcing all three errors to zero. A bounded
signal age preserves these gains. The exact scalar review flow has
$E^+-C=\alpha(E-C)$, $0<\alpha\le e^{-b_{\min}T_r}<1$, and therefore
the same tracking gain. Strict local small gain supplies stability as
well as attraction. S18 proves the comparison bounds, endpoint scope
and review argument in full. \hfill$\square$

At the common protective point $\Gamma\in(0.47499705,0.47499706)$.
Over the declared box $r\in[0.018,0.022]$, $K\in[90,110]$,
$q\in[0.0009,0.0011]$, $E^*\in[1.8,2.4]$,
$\delta\in[0.06,0.08]$, $M\in[26,34]$, $Z_{\rm ref}=1$,
with $E_0=E^*(1+\delta)$, outward interval arithmetic gives
$\Gamma<0.686294$. Thus neither a periodic attractor nor another
non-equilibrium limit set can compete with the equilibrium inside this
basin. This is a global protective result, not a global guarantee for
mobilising or mixed laws, Euler commands, or unbounded review gaps.
No convergence rate is uniform as delay/cadence diverges or initial
states approach the excluded faces.
% END BASIN

% BLOCK FOLD
\subsubsection{8.7 Infinite-series fold residuals and discretisation accuracy}
\label{infinite-fold-residuals}
An infinite-series residual estimate distinguishes finite-system
precision from continuum approximation error. Augmented orbit-and-kernel solves on odd grids
of 65, 129 and 257 points remove the even-grid Nyquist ambiguity and
refine both folds. At 257 points the approximations are
$\tau=5.587236198663494$, $T=315.322196288$ yr and
$\tau=64.40232665082372$, $T=72.885536774$ yr.
For the upper fold the off-grid orbit defect falls from
$1.30\times10^{-5}$ to $8.12\times10^{-9}$ to $6.52\times10^{-13}$
in physical time; the corresponding kernel defects are
$3.26\times10^{-4}$, $4.21\times10^{-7}$ and $2.32\times10^{-12}$.
These are convergence diagnostics, not existence certificates.

An interval Fourier calculation additionally bounds the \emph{whole
infinite series}, rather than only the sampled defect. Exact decimal
trigonometric coefficients define two approximate orbit/kernel pairs.
Analytic strips of half-width $0.008$ and $0.006$ exclude softplus branch
singularities and the rational pole. A 4096-point interval FFT, with
analytic alias and infinite-tail estimates, gives the following upper
bounds in the phase-variable weighted Wiener norm
$\|u\|_{1,\nu}=\sum_k|u_k|\nu^{|k|}$, $\nu=1.001$:
\begin{center}\small
\begin{tabular}{lrr}
\toprule Fold & largest orbit component & largest kernel component\\\midrule
lower & $3.170\times10^{-10}$ & $6.360\times10^{-12}$\\
upper & $2.560\times10^{-10}$ & $3.140\times10^{-12}$\\\bottomrule
\end{tabular}
\end{center}
These bounds include every Fourier mode of the nonlinear orbit and
linearised fold-kernel residuals. S19 gives the exact residual equations,
analytic inequalities, interval algorithm and phase constraints.

A continuum fold certificate still requires a validated inverse for the
full bordered infinite-dimensional operator, controlled finite-tail
coupling, a nonlinear contraction, and continuum transversality and
curvature. Delay shifts with free period/delay require derivative-loss
control; a naive finite Jacobian inverse is not that proof. The upper
257-point delay differs from the $m=64$ finite certificate by about
$5.5\times10^{-7}$ yr, much more than that certificate's interval width.
The narrow finite interval therefore measures finite-system accuracy,
not continuum discretisation error. No continuum existence or full
periodic Floquet-spectrum claim is made from these residuals.
% END FOLD

% BLOCK FIELD
\paragraph{Quantitative dynamic field tests.}
The same herring stock also permits a retrospective predictive test,
separate from event-clock interpretation. ICES SAG assessment 2025,
key 19373, supplies annual total biomass, catch and recruitment through
2024; the incomplete 2025 terminal row is excluded (ICES, 2025).
Because the total-biomass unit field is blank, the observation model is
explicitly dimensionless:
\[
 N_y=B_y/B_{\rm med},\qquad
 U_y=(C_y/C_{\rm med})/(B_y/B_{\rm med}).
\]
The reference medians use 1950--1994. $U$ is a catch-intensity proxy,
not measured institutional effort or age-2--6 fishing mortality.
Memory is driven by the observed decline $N_{y-1}-N_y$. The response
component predicts $U_{y+1}$ from origin-known $U_y$ and delayed memory
using an accurately integrated one-year scalar flow.

The common target split is 1960--1994 training, 1995--2009 validation
and 2010--2024 testing. Each law has two fitted coefficients and 260
memory/lag/ceiling candidates selected by validation, followed by a
pretest refit. Candidate-dependent missing early predictors in an initial
run prompted a disclosed common-support correction after those outputs
were inspected; both runs and the amendment are retained in S20's
protocol archive. This is a retrospective, non-preregistered study,
not a confirmatory test. No candidate is selected by test performance.
\begin{center}\small
\begin{tabular}{lrr}
\toprule Test endpoint, 2010--2024 & RMSE/persistence & RMSE/autoregression\\\midrule
mobilising response, one step & 0.900 & 1.035\\
protective response, one step & 0.865 & 0.995\\
logistic stock, one step & 1.355 & 1.355\\
recruitment-augmented stock, one step & 1.004 & 1.003\\
mobilising stock, recursive & 2.355 & 3.166\\
protective stock, recursive & 1.935 & 2.601\\\bottomrule
\end{tabular}
\end{center}
One-step autoregression is affine AR(1); recursive autoregression is a
two-state VAR(1), with constant-2009 persistence for the recursive
comparison. Endpoints and baselines must not be interchanged.

The protective response improves on persistence by about 13.5\%, but
beats AR(1) by only 0.47\% on 15 annual targets. The mobilising fit
hits its positive additive-coefficient lower bound and fails the AR(1)
comparison. Low/high biomass-path stress tests and time-step halving are
reported in S20; neither supplies independent replication. Secondary
exploratory tests use exact held-input stock flows
$\dot N=(r-hU)N-bN^2$, with and without lagged recruitment. Recruitment
reduces error, but neither ecological model beats both baselines.
Recursive annual sampled surrogates, initialised in 2009 without future
stock/catch observations, also fail the stock comparisons (Figure~\ref{fig:field-tests-v43}).
The protective recursive proxy alone beats VAR(1), which does not repair
its stock failure. Its fitted sufficient basin gain is about $2.31>1$,
outside Theorem 7.5; this does not prove instability.

Thus there is limited response-component predictive support, not joint
dynamic field validation of the autonomous DDE. Retrospective assessment
smoothing, the derived pressure proxy, recruitment and the absence of
direct institutional-state measurements remain material identification
obstacles. Timing compatibility, conditional prediction and joint mechanistic
validation therefore supply distinct evidence; the model-class theorem
does not itself validate a field mechanism.
\begin{figure}[htbp]
\centering\includegraphics[width=\linewidth]{fig4_field_tests_v43.png}
\caption{Blocked herring tests and exploratory recursive annual-surrogate
predictions. The top row uses observed prediction-origin states; the
bottom row starts in 2009 and uses no future stock/catch observations.
Biomass and catch intensity are separately normalised. Stock forecast
failure is retained even where the protective proxy prediction is
competitive. Curves are assessment-based retrospective tests, not
real-time historical forecasts or a fit of the continuous DDE.}
\label{fig:field-tests-v43}
\end{figure}
% END FIELD
